% GENERATED FILE. Edit canonical lessons, then run npm run generate:books.
% canonical-source-hash: fnv1a64:2fe4e5e50e493993
% canonical-lessons: ZH-C256-renwei, ZH-C256-kanfa, ZH-C256-kanlai, ZH-C256-yexu, ZH-C256-butong

\chapter{What Do You Think?}
\label{ch:zh-what-do-you-think}
\hlchaptermodality{256}

\begin{chapteropening}
\textbf{By the end of this chapter:} \emph{I can ask what someone thinks, give my view, say how things look, say perhaps, and say we are different.}

Complete the last lesson of chapter 256: I can ask what someone thinks, give my view, say how things look, say perhaps, and say we are different.
\end{chapteropening}

\section[rènwéi]{rènwéi — to think; to hold a view}

\label{lesson:ZH-C256-renwei}



\begin{quote}
Before the new one: say the Mandarin for to work overtime, then the Mandarin for the boss.
\end{quote}

\subsection*{You'll want to know: rènwéi}
\emph{rènwéi} — \textquotedblleft{}to think; to hold a view\textquotedblright{}.

\begin{quote}
\emph{nǐ rènwéi zěnmeyàng} — What do you think?
\end{quote}

Its characters come later; for now it is a word for the ear.

\begin{cousinweb}
\emph{rènwéi} is built from \emph{rèn} 'to know' and \emph{wéi}, both inherited from Middle Chinese, \emph{nyinH} and \emph{hjwe} (Guangyun, in Baxter's transcription), which glosses the first 'to know' and the second 'to make'.
\end{cousinweb}

\begin{grammarlens}[title={What you've built}]
One more word for saying no and asking.
\end{grammarlens}

\subsection*{Your turn}
\begin{itemize}
\raggedright
  \item \emph{Say it:} \emph{rènwéi}
  \item \emph{Say it:} \emph{rènwéi}, once more
  \item \emph{Say it:} \emph{nǐ rènwéi zěnmeyàng} — What do you think?
  \item \emph{Recall:} say the Mandarin for to work overtime, then the Mandarin for the boss, then say \emph{rènwéi} again
\end{itemize}

\begin{tcolorbox}[breakable,colback=teal!4,colframe=teal!35!black,title={Before you move on}]
What does rènwéi mean? (\textquotedblleft{}To think; to hold a view\textquotedblright{}.) Say it once more.
\end{tcolorbox}
\section[kànfǎ]{kànfǎ — view; opinion}

\label{lesson:ZH-C256-kanfa}



\begin{quote}
Before the new one: say the Mandarin for the boss, then the Mandarin for to hold a view.
\end{quote}

\subsection*{You'll want to know: kànfǎ}
\emph{kànfǎ} — \textquotedblleft{}view; opinion\textquotedblright{}.

\begin{quote}
\emph{wǒmen de kànfǎ yíyàng} — We see it the same way.
\end{quote}

Its characters come later; for now it is a word for the ear.

\begin{cousinweb}
\emph{kànfǎ} is built from \zh{看} 'to look' and \emph{fǎ} 'a rule, a way', both inherited from Middle Chinese, \emph{khanH} and \emph{pjop} (Guangyun, in Baxter's transcription), which glosses the second 'a rule': a way of looking.
\end{cousinweb}

\begin{grammarlens}[title={What you've built}]
One more word for saying no and asking.
\end{grammarlens}

\subsection*{Your turn}
\begin{itemize}
\raggedright
  \item \emph{Say it:} \emph{kànfǎ}
  \item \emph{Say it:} \emph{kànfǎ}, once more
  \item \emph{Say it:} \emph{wǒmen de kànfǎ yíyàng} — We see it the same way.
  \item \emph{Recall:} say the Mandarin for the boss, then the Mandarin for to hold a view, then say \emph{kànfǎ} again
\end{itemize}

\begin{tcolorbox}[breakable,colback=teal!4,colframe=teal!35!black,title={Before you move on}]
What does kànfǎ mean? (\textquotedblleft{}View; opinion\textquotedblright{}.) Say it once more.
\end{tcolorbox}
\section[\texorpdfstring{\zh{看来} (kànlai)}{看来 (kànlai)}]{\zh{看来} (kànlai) — it seems; it looks as if}

\label{lesson:ZH-C256-kanlai}



\begin{quote}
Before the new one: say the Mandarin for to hold a view, then the Mandarin for a point of view.
\end{quote}

\subsection*{You'll want to know: \zh{看来}}
\textbf{\zh{看来}} — \emph{kànlai} — \textquotedblleft{}it seems; it looks as if\textquotedblright{}.

\begin{quote}
\textbf{\zh{看来他今天不来了}} \emph{kànlai tā jīntiān bù lái le} — It looks as if he isn't coming today.
\end{quote}

\begin{cousinweb}
\textbf{\zh{看来}} (\emph{kànlai}) is built from \zh{看} 'to look' and \zh{来} 'to come', both inherited from Middle Chinese, \emph{khanH} and \emph{loj} (Guangyun, in Baxter's transcription). In the word the second loses its tone.
\end{cousinweb}

\begin{grammarlens}[title={What you've built}]
One more word for saying no and asking.
\end{grammarlens}

\subsection*{Your turn}
\begin{itemize}
\raggedright
  \item \emph{Say it:} \emph{kànlai}
  \item \emph{Say it:} \emph{kànlai}, once more
  \item \emph{Say it:} \emph{kànlai tā jīntiān bù lái le} — It looks as if he isn't coming today.
  \item \emph{Recall:} say the Mandarin for to hold a view, then the Mandarin for a point of view, then say \emph{kànlai} again
\end{itemize}

\begin{tcolorbox}[breakable,colback=teal!4,colframe=teal!35!black,title={Before you move on}]
What does \zh{看来} mean? (\textquotedblleft{}It seems; it looks as if\textquotedblright{}.) Say it once more.
\end{tcolorbox}
\section[yěxǔ]{yěxǔ — it may be that; perhaps}

\label{lesson:ZH-C256-yexu}



\begin{quote}
Before the new one: say the Mandarin for a point of view, then the Mandarin for it looks as if.
\end{quote}

\subsection*{You'll want to know: yěxǔ}
\emph{yěxǔ} — \textquotedblleft{}it may be that; perhaps\textquotedblright{}.

\begin{quote}
\emph{yěxǔ tā wàng le} — Maybe he forgot.
\end{quote}

Its characters come later; for now it is a word for the ear.

Two third tones meet in this word: yě is said with a rising tone before xǔ, and the spelling keeps both marks.

\begin{cousinweb}
\emph{yěxǔ} is built from \zh{也} 'also' and \emph{xǔ} 'to allow', both inherited from Middle Chinese, \emph{yaeX} and \emph{xjoX} (Guangyun, in Baxter's transcription), which glosses the first as a particle that closes a sentence and the second 'to permit'. How the pair came to mean perhaps is not settled.
\end{cousinweb}

\begin{grammarlens}[title={What you've built}]
One more word for saying no and asking.
\end{grammarlens}

\subsection*{Your turn}
\begin{itemize}
\raggedright
  \item \emph{Say it:} \emph{yěxǔ}
  \item \emph{Say it:} \emph{yěxǔ}, once more
  \item \emph{Say it:} \emph{yěxǔ tā wàng le} — Maybe he forgot.
  \item \emph{Recall:} say the Mandarin for a point of view, then the Mandarin for it looks as if, then say \emph{yěxǔ} again
\end{itemize}

\begin{tcolorbox}[breakable,colback=teal!4,colframe=teal!35!black,title={Before you move on}]
What does yěxǔ mean? (\textquotedblleft{}It may be that; perhaps\textquotedblright{}.) Say it once more.
\end{tcolorbox}
\section[\texorpdfstring{\zh{不同} (bùtóng)}{不同 (bùtóng)}]{\zh{不同} (bùtóng) — different}

\label{lesson:ZH-C256-butong}



\begin{quote}
Before the new one: say the Mandarin for it looks as if, then the Mandarin for it may be that.
\end{quote}

\subsection*{You'll want to know: \zh{不同}}
\textbf{\zh{不同}} — \emph{bùtóng} — \textquotedblleft{}different\textquotedblright{}.

\begin{quote}
\textbf{\zh{我们两个人很不同}} \emph{wǒmen liǎng ge rén hěn bùtóng} — The two of us are very different.
\end{quote}

Tone: \zh{不} keeps its own falling tone here, bù, because tóng after it is a rising tone, not a falling one.

\begin{cousinweb}
\textbf{\zh{不同}} (\emph{bùtóng}) is built from \zh{不} 'not' and \zh{同} 'the same', both inherited from Middle Chinese, \emph{pjut} and \emph{duwng} (Guangyun, in Baxter's transcription).
\end{cousinweb}

\begin{grammarlens}[title={What you've built}]
One more word for saying no and asking.
\end{grammarlens}

\subsection*{Your turn}
\begin{itemize}
\raggedright
  \item \emph{Say it:} \emph{bùtóng}
  \item \emph{Say it:} \emph{bùtóng}, once more
  \item \emph{Say it:} \emph{wǒmen liǎng ge rén hěn bùtóng} — The two of us are very different.
  \item \emph{Recall:} say the Mandarin for it looks as if, then the Mandarin for it may be that, then say \emph{bùtóng} again
\end{itemize}

\begin{tcolorbox}[breakable,colback=teal!4,colframe=teal!35!black,title={Before you move on}]
What does \zh{不同} mean? (\textquotedblleft{}Different\textquotedblright{}.) Say it once more.
\end{tcolorbox}
